% =====================================================================
%  Master's Thesis Template — Applied Machine Learning
%  Supervisor: Alex Jung, Aalto University
%
%  This skeleton is structured around the SCI grade characterisation
%  (see material/GradeCharact.pdf) and the writing conventions in
%  README.md. Inline TODO comments map each section to the grading
%  dimensions:
%      [G1] Definition of research scope and goals
%      [G2] Command of the topic
%      [G3] Methods
%      [G4] Results and contribution
%      [G5] Presentation, language and structure
%
%  Replace every TODO with your own work. Compile with:
%      pdflatex thesis && bibtex thesis && pdflatex thesis && pdflatex thesis
% =====================================================================

\documentclass[11pt,a4paper,oneside]{report}

% --- Encoding & fonts ---
\usepackage[utf8]{inputenc}
\usepackage[T1]{fontenc}
\usepackage{lmodern}

% --- Layout ---
\usepackage[a4paper,margin=2.5cm]{geometry}
\usepackage{setspace}
\onehalfspacing

% --- Math & symbols ---
\usepackage{amsmath,amssymb,amsthm}
\usepackage{mathtools}

% --- Floats, figures, tables ---
\usepackage{graphicx}
\usepackage{caption}
\usepackage{subcaption}
\usepackage{booktabs}

% --- Algorithms (for pseudocode) ---
\usepackage{algorithm}
\usepackage{algpseudocode}

% --- Bibliography (IEEE style: see README "References" section) ---
\usepackage[backend=biber,style=ieee,sorting=none]{biblatex}
\addbibresource{references.bib}    % create this file alongside the .tex

% --- Hyperlinks ---
\usepackage[hidelinks]{hyperref}
\usepackage{cleveref}

% --- Convenience macros ---
\newcommand{\featurevec}{\mathbf{x}}      % feature vector
\newcommand{\labelscalar}{y}              % label
\newcommand{\dataset}{\mathcal{D}}        % training set
\newcommand{\hypothesis}{h}               % learned predictor
\newcommand{\loss}{\ell}                  % loss function
\newcommand{\reals}{\mathbb{R}}
\newcommand{\expect}{\mathbb{E}}

\begin{document}

% =====================================================================
%  TITLE PAGE — replace placeholders with your details
% =====================================================================
\begin{titlepage}
\centering
\vspace*{2cm}
{\LARGE\bfseries TODO: Thesis Title \\[0.4em]
\large An informative subtitle that names the problem and the method\par}
\vspace{2cm}
{\Large TODO: Author Name\par}
\vspace{0.3cm}
{\large Master's Thesis\par}
{\large Aalto University, School of Science\par}
{\large Master's Programme in Computer, Communication and Information Sciences\par}
\vfill
{\large Supervisor: Prof.\ Alex Jung\par}
{\large Advisor: TODO\par}
\vspace{0.5cm}
{\large \today\par}
\end{titlepage}

% =====================================================================
%  ABSTRACT
%  ~250 words. State (i) the problem, (ii) the data, (iii) the method,
%  (iv) the headline result, (v) the takeaway. No citations, no
%  acronyms without expansion.
%  [G1, G4]
% =====================================================================
\chapter*{Abstract}
\addcontentsline{toc}{chapter}{Abstract}

% TODO [G1]: state the ML problem in one sentence.
% TODO [G1]: state the data source and what data points represent.
% TODO [G3]: name the method(s) and any baselines.
% TODO [G4]: report the headline result (one number, with units / metric).
% TODO [G4]: state the practical or scientific implication in one sentence.

% =====================================================================
%  AI-USE DISCLOSURE
%  Goes here as a dedicated statement — NOT inside Methods.
%  See README "Responsible Use of AI".
% =====================================================================
\chapter*{Statement on the Use of AI Tools}
\addcontentsline{toc}{chapter}{Statement on the Use of AI Tools}

% TODO: list each AI tool used (name, version, settings, purpose).
% Example template:
%
% During the preparation of this thesis the author used the following
% AI tools:
%   - <Tool name, version> (<provider>): used for <purpose, e.g. proofreading
%     of language in Chapters 2 and 4; brainstorming literature search terms;
%     generating boilerplate plotting code for Figures 3.2 and 4.1>.
%
% All AI-generated content was independently reviewed, verified, and
% revised by the author. The author bears full responsibility for the
% content of this thesis.

% =====================================================================
%  ACKNOWLEDGEMENTS (optional)
% =====================================================================
\chapter*{Acknowledgements}
\addcontentsline{toc}{chapter}{Acknowledgements}

% TODO: thank your supervisor, advisor, family, etc.

% =====================================================================
%  TABLE OF CONTENTS
% =====================================================================
\tableofcontents

% =====================================================================
%  CHAPTER 1 — INTRODUCTION
%  Open with a paragraph that sets the stage. Then: motivation,
%  research goals, contributions, outline.  [G1]
% =====================================================================
\chapter{Introduction}
\label{ch:intro}

% Section opener: in 2--4 sentences, what is this chapter about and
% how does it connect to the rest of the thesis?  [G5]

\section{Motivation}
% TODO [G1]: why does this problem matter? Use one concrete example
% (a real-world setting, a stakeholder, or a measurable cost of failure).

\section{Research Goals and Questions}
\label{sec:goals}
% TODO [G1]: state 1--3 research questions explicitly. A grade-5 thesis
% has "precisely defined and justified research questions". Avoid hedged
% phrasing like "we explore". Prefer "we investigate whether X improves Y
% relative to baseline Z under conditions W".
%
% Example:
%   \textbf{RQ1.} Does method $M$ achieve lower test error than baseline
%   $B$ on dataset $D$ under a fixed compute budget?
%   \textbf{RQ2.} Which features in $D$ are most predictive of $\labelscalar$,
%   and is this consistent with prior findings in [CITE]?

\section{Contributions}
% TODO [G4]: list 2--4 concrete contributions. Each should be verifiable
% from the thesis body. Examples:
%   1. An empirical comparison of methods $M_1, M_2, M_3$ on dataset $D$
%      against baselines $B_1, B_2$ (Chapter~\ref{ch:results}).
%   2. A diagnostic analysis showing that error grows with $\ldots$
%      (Section~\ref{sec:diagnosis}).
%   3. An open-source implementation: \texttt{github.com/<you>/<repo>}.

\section{Thesis Outline}
% Brief paragraph: what each chapter covers.

% =====================================================================
%  CHAPTER 2 — BACKGROUND AND RELATED WORK
%  Demonstrates "command of the topic". [G2]
%  Cite primarily peer-reviewed venues; use IEEE format.
% =====================================================================
\chapter{Background and Related Work}
\label{ch:background}

% Opener: what does the reader need to know before Chapter 3, and what
% prior work does this thesis build on or contrast with?

\section{Preliminaries}
% TODO [G2]: define notation and the ML concepts you will use.
% Use terminology from the Aalto Dictionary of ML
% (https://aaltodictionaryofml.github.io/).

\section{Related Work}
% TODO [G2]: organise by theme, not chronologically. End each
% subsection with a sentence stating the gap your thesis addresses.
% A grade-5 thesis "throws light on the topic from a variety of perspectives"
% — show that you know the alternative framings, not only the one you chose.

% =====================================================================
%  CHAPTER 3 — PROBLEM FORMULATION
%  Define the ML problem precisely. README explicitly asks for this.
%  [G1, G3]
% =====================================================================
\chapter{Problem Formulation}
\label{ch:problem}

% Opener.

\section{Data Points, Features, and Labels}
% TODO [G1]: state precisely what one data point is, what its features
% are, and what the label is. This is non-negotiable. Example:
%
% A data point is a single hourly weather observation at station $s$ and
% hour $t$. Its feature vector $\featurevec \in \reals^d$ stacks the past
% 24 hourly readings of temperature, pressure, and wind speed
% ($d = 72$). The label $\labelscalar \in \reals$ is the temperature 6
% hours ahead.

\section{Loss Functions}
% TODO [G1, G3]: explicitly name the loss used for training and,
% separately, the metric used for validation/testing.
%
% Example:
% Training minimises the empirical risk
% \begin{equation}
% \label{eq:train-loss}
% \widehat{L}(\hypothesis) = \frac{1}{|\dataset|} \sum_{(\featurevec,\labelscalar) \in \dataset}
%   \loss\!\big(\hypothesis(\featurevec),\, \labelscalar\big),
% \end{equation}
% with the squared loss $\loss(\hat y, y) = (\hat y - y)^2$. Validation
% and test performance are reported using the root-mean-squared error
% (RMSE), not the squared loss, to match prior work [CITE].
%
% Note the punctuation: \eqref{eq:train-loss} ends in a comma because
% the surrounding sentence continues; use a period if it ends the sentence.
% Number an equation only if you reference it in the text.

\section{Assumptions and Scope}
% TODO [G1]: state assumptions (i.i.d.\ data? stationarity? labelled
% data? compute budget?) and what is out of scope.

% =====================================================================
%  CHAPTER 4 — METHODS
%  Describe what you did, not why ML is interesting. [G3]
%  Pseudocode for any new algorithm.
% =====================================================================
\chapter{Methods}
\label{ch:methods}

% Opener.

\section{Method Overview}
% TODO [G3]: 1--2 paragraphs naming the methods used and why each is
% appropriate for the problem in Chapter \ref{ch:problem}.

\section{Algorithm}
\label{sec:algo}
% TODO [G3]: present any new algorithm as pseudocode.
% See https://www.overleaf.com/learn/latex/Algorithms

\begin{algorithm}[h]
\caption{TODO: Name of your algorithm}
\label{alg:main}
\begin{algorithmic}[1]
\Require Training set $\dataset$, learning rate $\eta$, iterations $T$
\Ensure  Trained predictor $\hypothesis$
\State Initialise $\hypothesis^{(0)}$
\For{$t = 1, \dots, T$}
    \State Sample mini-batch $\mathcal{B} \subset \dataset$
    \State Compute $g_t \gets \nabla_{\hypothesis}\, \widehat{L}(\hypothesis^{(t-1)}; \mathcal{B})$
    \State $\hypothesis^{(t)} \gets \hypothesis^{(t-1)} - \eta\, g_t$
\EndFor
\State \Return $\hypothesis^{(T)}$
\end{algorithmic}
\end{algorithm}

\section{Implementation}
% TODO: framework, libraries, hardware. Enough detail that someone
% else could re-run your experiments. Link to the public code repo.

% =====================================================================
%  CHAPTER 5 — EXPERIMENTAL SETUP
%  Data, baselines, metrics, train/val/test split, hyperparameters.
%  [G3, G4]
% =====================================================================
\chapter{Experimental Setup}
\label{ch:setup}

% Opener.

\section{Dataset}
% TODO [G4]: source, size, splits, preprocessing, label distribution.
% Show one or two figures: a histogram of labels, a few example data points.

\section{Baselines}
% TODO [G4]: README emphasises baselines. Use at least:
%   (i)  a trivial baseline (mean predictor, majority class, or persistence),
%   (ii) a strong but simple baseline (linear/logistic regression, gradient boosting),
%   (iii) prior published results on the same task, if any.
% Justify each choice in one sentence.

\section{Evaluation Metrics}
% TODO [G4]: name the metric, define it (with an equation if non-standard),
% and say why it matches the goal in Chapter \ref{ch:problem}.

\section{Hyperparameter Selection}
% TODO: how were hyperparameters tuned? On which split? Report the search
% space, not just the final values.

% =====================================================================
%  CHAPTER 6 — RESULTS
%  Numbers, tables, plots — and discussion of what they mean. [G4]
% =====================================================================
\chapter{Results}
\label{ch:results}

% Opener.

\section{Headline Results}
% TODO [G4]: a single table comparing your method against all baselines
% on the test set, with confidence intervals or std over seeds. Discuss
% what is and is not statistically meaningful.

% Example skeleton (uncomment and fill):
% \begin{table}[h]
% \centering
% \caption{Test-set performance on dataset $D$ (mean $\pm$ std over 5 seeds).
%          Lower is better.}
% \label{tab:headline}
% \begin{tabular}{lcc}
% \toprule
% Method     & Metric A $\downarrow$ & Metric B $\downarrow$ \\
% \midrule
% Trivial    & $\ldots$ & $\ldots$ \\
% Baseline   & $\ldots$ & $\ldots$ \\
% Ours       & $\ldots$ & $\ldots$ \\
% \bottomrule
% \end{tabular}
% \end{table}

\section{Diagnostic Analysis}
\label{sec:diagnosis}
% TODO [G4]: where does the model fail? Break down errors by subgroup,
% input regime, or class. README calls for "model diagnosis using
% numerical experiments (benchmarks, sensitivity analysis) and, when
% appropriate, mathematical analysis".

\section{Sensitivity and Ablations}
% TODO [G4]: vary one design choice at a time (data size, feature subset,
% hyperparameter) and plot the effect.

% =====================================================================
%  CHAPTER 7 — DISCUSSION
%  Place results in context. [G2, G4]
% =====================================================================
\chapter{Discussion}
\label{ch:discussion}

\section{Answers to the Research Questions}
% TODO: one paragraph per RQ from Section \ref{sec:goals}, answering it
% with a concrete pointer into Chapter \ref{ch:results}.

\section{Comparison to Prior Work}
% TODO [G2]: how do your results compare to those cited in
% Chapter \ref{ch:background}? If they differ, explain why.

\section{Limitations}
% TODO: be specific. "More data would help" is not a limitation;
% "the dataset covers only winter months, so generalisation to summer
% is untested" is.

\section{Threats to Validity}
% TODO: data leakage checks, label noise, distribution shift between
% train and test, possible confounders.

% =====================================================================
%  CHAPTER 8 — CONCLUSION
% =====================================================================
\chapter{Conclusion}
\label{ch:conclusion}

\section{Summary}
% TODO: 1--2 paragraphs. What was the problem, what did you do, what did
% you find?

\section{Future Work}
% TODO: 2--3 concrete next steps. "Try transformers" is not concrete;
% "extend the comparison to time-series transformer architectures
% (e.g., [CITE]) on the same dataset" is.

% =====================================================================
%  REFERENCES (IEEE format)
%  Maintain a separate references.bib file.
% =====================================================================
\printbibliography[heading=bibintoc,title={References}]

% =====================================================================
%  APPENDICES (optional)
%  Use for derivations, full hyperparameter grids, extra plots.
% =====================================================================
\appendix
\chapter{Additional Experimental Details}
% TODO

\end{document}